\section*{Exercises on \S\arabic{exercisegroup}}
\phantomsection
\addcontentsline{toc}{section}{Exercises on \protect\S\arabic{exercisegroup}}

\begin{enumerate}
\def\labelenumi{\arabic{enumi})}
\tightlist
\item
  Let \(f\) be a vector function of a real variable, defined on an interval \(I \subset \mathbf{R}\) and differentiable at a point \(x_0\) interior to \(I\) . Show that the quotient
\end{enumerate}

\[
\frac {f (x _ {0} + h) - f (x _ {0} - k)}{h + k}
\]

tends to \(f'(x_0)\) as \(h\) and \(k\) tend to 0 through values \(> 0\) . Converse.

*Show that the function \(f\) equal to \(x^2 \sin 1/x\) for \(x \neq 0\) , and to 0 for \(x = 0\) , is everywhere differentiable, but that \((f(y) - f(z)) / (y - z)\) does not approach \(f'(0)\) as \(y\) and \(z\) tend to 0, while remaining distinct and \(>0_*\)

\begin{enumerate}
\def\labelenumi{\arabic{enumi})}
\setcounter{enumi}{1}
\tightlist
\item
  On the interval I = {[}0, 1{]} we define a sequence of continuous real functions \((f_{n})\) inductively as follows: We take \(f_{0}(x) = x\) ; for each integer \(n \geqslant 1\) the function \(f_{n}\) is affine linear on each of the \(3^{n}\) intervals \(\left[\frac{k}{3^{n}}, \frac{k+1}{3^{n}}\right]\) for \(0 \leqslant k \leqslant 3^{n}-1\) ; further, we take
\end{enumerate}

\[
f _ {n + 1} \left(\frac {k}{3 ^ {n}}\right) = f _ {n} \left(\frac {k}{3 ^ {n}}\right)
\]

\[
f _ {n + 1} \left(\frac {k}{3 ^ {n}} + \frac {1}{3 ^ {n + 1}}\right) = f _ {n} \left(\frac {k}{3 ^ {n}} + \frac {2}{3 ^ {n + 1}}\right), f _ {n + 1} \left(\frac {k}{3 ^ {n}} + \frac {2}{3 ^ {n + 1}}\right) = f _ {n} \left(\frac {k}{3 ^ {n}} + \frac {1}{3 ^ {n + 1}}\right).
\]

Show that the sequence \((f_{n})\) converges uniformly on I to a continuous function which has no derivative (finite or infinite) at any point of the interval {]}0, 1{[} (use exerc. 1).

\begin{enumerate}
\def\labelenumi{\arabic{enumi})}
\setcounter{enumi}{2}
\item
  Let \(\mathcal{C}(\mathrm{I})\) be the complete space of continuous finite real functions defined on the compact interval \(\mathrm{I} = [a,b]\) of \(\mathbf{R}\) , and endow \(\mathcal{C}(\mathrm{I})\) with the topology of uniform convergence (Gen.~Top., X, p.~277). Let A be the subset of \(\mathcal{C}(\mathrm{I})\) formed by the functions \(x\) such that for at least one point \(t\in [a,b[\) (depending on the function \(x\) ) the function \(x\) has a finite right derivative. Show that A is a meagre set in \(\mathcal{C}(\mathrm{I})\) (Gen.~Top., IX, p.~192), and hence its complement, that is, the set of continuous functions on I not having a finite right derivative at any point of \([a,b[\) is a Baire subspace of \(\mathcal{C}(\mathrm{I})\) (Gen.~Top., IX, p.~192). (Let \(\mathrm{A}_n\) be the set of functions \(x\in \mathcal{C}(\mathrm{I})\) such that for at least one value of \(t\) satisfying \(a\leqslant t\leqslant b - 1 / n\) (and depending on \(x\) ) one has \(|x(t') - x(t)|\leqslant n|t' - t|\) for all \(t'\) such that \(t\leqslant t'\leqslant t + 1 / n\) . Show that each \(\mathrm{A}_n\) is a closed nowhere dense set in \(\mathcal{C}(\mathrm{I})\) : remark that in \(\mathcal{C}(\mathrm{I})\) each ball contains a function having bounded right derivative on \([a,b[\) ; on the other hand, for every \(\varepsilon >0\) and every integer \(m > 0\) there exists on I a continuous function having at every point of \([a,b[\) a finite right derivative such that, for all \(t\in [a,b[\) one has \(|y(t)|\leqslant \varepsilon\) and \(|y_r'(t)|\geqslant m\) .)
\item
  Let E be a topological vector space over R and f a continuous vector function defined on an open interval \(I \subset R\) , and having a right derivative and a left derivative at every point of I.
\end{enumerate}

\begin{enumerate}
\def\labelenumi{\alph{enumi})}
\item
  Let U be a nonempty open set in E, and A the subset of I formed by the points x such that \(\mathbf{f}_{d}^{\prime}(x) \in \mathrm{U}\) . Given a number \(\alpha > 0\) let B be the subset of I formed by the points x such that there exists at least one \(y \in I\) satisfying the conditions \(x - \alpha \leqslant y < x\) and \((\mathbf{f}(x) - \mathbf{f}(y)) / (x - y) \in \mathrm{U}\) ; show that the set B is open and that \(A \cap CB\) is countable (remark that this last set is formed by the left-hand endpoints of intervals contiguous to CB). Deduce that the set of points \(x \in A\) such that \(\mathbf{f}_{g}^{\prime}(x) \notin \overline{\mathbf{U}}\) is countable.
\item
  Suppose that E is a normed space; the image \(\mathbf{f}(\mathrm{I})\) is then a metric space having a countable base, and the same is true for the closed vector subspace F of E generated by \(\mathbf{f}(\mathrm{I})\) , a subspace which contains \(\mathbf{f}_{d}^{\prime}(\mathrm{I})\) and \(\mathbf{f}_{g}^{\prime}(\mathrm{I})\) . Deduce from a) that the set of points \(x \in I\) such that \(\mathbf{f}_{d}^{\prime}(x) \neq \mathbf{f}_{g}^{\prime}(x)\) is countable. (If \((\mathrm{U}_{m})\) is a countable base for the topology of F note that for two distinct points a, b of F there exist two disjoint sets \(U_{p}, U_{q}\) such that \(a \in U_{p}\) and \(b \in U_{q}\) .)
\item
  Take for E the product \(R^{I}\) (the space of mappings from I into R, endowed with the topology of simple convergence), and for each \(x \in I\) denote by \(\mathbf{g}(x)\) the map \(t \mapsto |x - t|\) of I into R. Show that g is continuous and that, for every \(x \in I\) , one has \(\mathbf{g}_{r}'(x) \neq \mathbf{g}_{l}'(x)\) .
\end{enumerate}

\begin{enumerate}
\def\labelenumi{\arabic{enumi})}
\setcounter{enumi}{4}
\tightlist
\item
  Let \(\mathbf{f}\) be a continuous vector function defined on an open interval \(\mathrm{I} \subset \mathbf{R}\) with values in a normed space \(\mathrm{E}\) over \(\mathbf{R}\) , and admitting a right derivative at every point of \(\mathrm{I}\) .
\end{enumerate}

\begin{enumerate}
\def\labelenumi{\alph{enumi})}
\item
  Show that the set of points \(x \in \mathrm{I}\) such that \(\mathbf{f}_d'\) is bounded on a neighbourhood of \(x\) is an open set dense in \(\mathrm{I}\) (use th. 2 of Gen.~Top., IX, p.~194).
\item
  Show that the set of points of I where \(\mathbf{f}_d^{\prime}\) is continuous is the complement of a meagre subset of I (cf.~Gen.~Top., IX, p.~255, exerc. 21).
\end{enumerate}

\begin{enumerate}
\def\labelenumi{\arabic{enumi})}
\setcounter{enumi}{5}
\tightlist
\item
  Let \((r_n)\) be the sequence formed by the rational numbers in {[}0, 1{]}, arranged in a certain order. Show that the function \(f(x) = \sum_{n=0}^{\infty} 2^{-n}(x - r_n)^{1/3}\) is continuous and differentiable
\end{enumerate}

at every point of \(\mathbf{R}\) , and has an infinite derivative at every point \(r_n\) . (To see that \(f\) is differentiable at a point \(x\) distinct from the \(r_n\) , distinguish two cases, according to whether the series with general term \(2^{-n}(x - r_n)^{-2/3}\) has sum \(+\infty\) or converges; in the second case, note for all \(x \neq 0\) and all \(y \neq x\) , one has

\[
0 \leqslant (y ^ {1 / 3} - x ^ {1 / 3}) / (y - x) \leqslant 4 / 3 x ^ {2 / 3}).
\]

\begin{enumerate}
\def\labelenumi{\arabic{enumi})}
\setcounter{enumi}{6}
\item
  Let \(f\) be a real function defined on an interval \(\mathrm{I} \subset \mathbf{R}\) , admitting a right derivative \(f_d'(x_0) = 0\) at a point of \(\mathrm{I}\) , and let \(\mathbf{g}\) be a vector function defined on a neighbourhood of \(y_0 = f(x_0)\) , having a right derivative and a left derivative (not necessarily equal) at this point. Show that \(\mathbf{g} \circ f\) has a right derivative equal to 0 at the point \(x_0\) .
\item
  Let \(f\) be a mapping from \(\mathbf{R}\) to itself such that the set \(C\) of points of \(\mathbf{R}\) where \(f\) is continuous is dense in \(\mathbf{R}\) , and such that the complement \(A\) of \(C\) is also dense. Show that the set \(D\) of points of \(C\) where \(f\) is right differentiable is meagre. (For each integer \(n\) , let \(E_n\) be the set of points \(a \in \mathbf{R}\) such that there exist two points \(x, y\) such that \(0 < x - a < 1/n\) , \(0 < y - a < 1/n\) and
\end{enumerate}

\[
\frac {f (x) - f (a)}{x - a} - \frac {f (y) - f (a)}{y - a} > 1.
\]

Show that the interior of \(E_{n}\) is dense in R. For this, note that for every open nonempty interval I in R there is a point \(b \in I \cap A\) ; show that for a \textless{} b and b - a sufficiently small one has \(a \in E_{n}\) .

\begin{enumerate}
\def\labelenumi{\arabic{enumi})}
\setcounter{enumi}{8}
\tightlist
\item
  Let \(\mathcal{B}(\mathbf{N})\) be the space of bounded sequences \(\mathbf{x} = (x_n)_{n\in \mathbb{N}}\) of real numbers, endowed with the norm \(\| \mathbf{x}\| = \sup |x_n|\) ; give an example of a continuous map \(t\mapsto \mathbf{f}(t) = (f_n(t))_{n\in \mathbb{N}}\) of \(\mathbf{R}\) into \(\mathcal{B}(\mathbf{N})\) such that each of the functions \(f_{n}\) is differentiable for \(t = 0\) , but \(\mathbf{f}\) is not differentiable at this point.
\end{enumerate}

\setcounter{exercisegroup}{1}
\refstepcounter{exercisegroup}
